% Part: lambda-calculus
% Chapter: introduction
% Section: lambda-definability

\documentclass[../../../include/open-logic-section]{subfiles}

\begin{document}

\olfileid{lam}{int}{rep}
\olsection{\usetoken{S}{lambda definable} અંકગણિતીય વિધેયો}

લૅમ્બડા કલન સંગણનના નિદર્શન તરીકે કેવી રીતે કામ કરી શકે? શરૂઆતમાં
આ વિધાનનો અર્થ શું એવો પ્રશ્ન પણ થાય. પ્રાકૃતિક સંખ્યાઓ પરની
સંગણનીયતા વિશે વાત કરવા માટે આ સંખ્યાઓનું યોગ્ય નિરૂપણ શોધવું
જોઈએ. નીચેનું નિરૂપણ આશ્ચર્યજનક રીતે સારું કામ કરે છે.

\begin{defn}
દરેક પ્રાકૃતિક સંખ્યા~$n$ માટે \emph{ચર્ચ અંક}
$\num{n}$ એ લૅમ્બડા પદ $\lambd[x][\lambd[y][(x(x(x(\dots
x(y)))))]]$ હોય એમ વ્યાખ્યાયિત કરો; તેમાં કુલ~$n$ વાર~$x$ આવે છે.
\end{defn}

$\num{n}$ પદો ``પુનરાવર્તકો'' છે: દલીલ~$f$ આપીએ તો~$\num{n}$
$y$ને~$f^n(y)$માં મોકલતું વિધેય આપે છે. દરેક ચર્ચ અંક સામાન્ય
સ્વરૂપમાં છે એ ધ્યાનમાં લો. હવે પ્રાકૃતિક સંખ્યાઓ પરના વિધેયને
લૅમ્બડા પદ ``ગણે છે'' તેનો અર્થ કહી શકીએ.

\begin{defn}
$f(x_0, \dots, x_{k-1})$ એ~$\Nat^k$માંથી~$\Nat$માં જતું
$k$-સ્થાની આંશિક વિધેય હોય. $\lambd$-પદ~$F$ એ~$f$ને
\emph{!!{lambda define}s} કરે છે એમ ત્યારે અને માત્ર ત્યારે જ
કહીએ જ્યારે પ્રાકૃતિક સંખ્યાઓની દરેક શ્રેણી $n_0$, \dots,~$n_{k-1}$
માટે,
\[
F\, \num{n_0}\, \num{n_1} \dots \num{n_{k-1}} \red \num{f(n_0, n_1, \dots,
  n_{k-1})}
\]
જો $f(n_0, \dots, n_{k-1})$ વ્યાખ્યાયિત હોય, અને ન હોય તો
$F\, \num{n_0}\, \num{n_1} \dots \num{n_{k-1}}$ને કોઈ સામાન્ય
સ્વરૂપ ન હોય.
\end{defn}
\sourcecorrection{OLLAM-003}{સ્થિર અંગ્રેજી મૂળ અહીં
$f(x_0,\dots,x_{k-1})$ને~$n$-સ્થાની કહે છે, પરંતુ દલીલોની
સંખ્યા~$k$ છે. આ કિસ્સામાં વિધેયનું ઇનપુટ ક્ષેત્ર~$\Nat^k$
છે. અહીં બંને બાબતો સ્પષ્ટ કરી છે.}
\sourcecorrection{OLLAM-004}{સ્થિર અંગ્રેજી મૂળમાં વિધેય
અવ્યાખ્યાયિત હોય ત્યારે~$F$ પછી અલ્પવિરામ મૂકીને દલીલોની યાદી
લખી છે. અહીં અલ્પવિરામ દૂર કરીને અગાઉની પંક્તિની જેમ જ
$F$નો ચર્ચ અંકો પરનો અનુપ્રયોગ લખ્યો છે.}

\begin{thm}
\ollabel{thm:lambda-def}
વિધેય~$f$ આંશિક સંગણનીય હોય ત્યારે અને માત્ર ત્યારે જ તે
લૅમ્બડા પદ વડે !!{lambda defined} હોય.
\end{thm}

\begin{explain}
આ પ્રમેય કંઈક અંશે ચોંકાવનારું છે. સંગણનના નિદર્શન તરીકે
લૅમ્બડા કલન ઘણું સરળ કલન છે: તેમાં માત્ર લૅમ્બડા અમૂર્તીકરણ અને
અનુપ્રયોગ જ ક્રિયાઓ છે!{} છતાં આટલાં અલ્પ સાધનોથી કોઈપણ
સંગણકીય પ્રક્રિયા અમલમાં મૂકી શકાય છે.
\end{explain}

\end{document}
